\documentclass{article}
\usepackage{graphicx} % Required for inserting images

\title{Modeling Real Travel Routes Using Shortest-Path Algorithms}
\author{Jennifer M.}
\date{February 2026}

\begin{document}

\maketitle

\section{Abstract}
Shortest-path algorithms are a fundamental topic in graph theory and are widely used to model routing problems in transportation networks. In these models, road systems are represented as weighted graphs, and routes are chosen to minimize quantities such as distance or travel time. While these algorithms are efficient and mathematically understood, the usefulness depends on how well their assumptions reflect real travel conditions. This paper examines how classical shortest-path algorithms are applied to transportation networks and discusses the limitations of these models, specifically the assumption that travel costs are fixed.  

\section{Introduction}
Graphs provide a natural way to represent networks, including transportation systems. Locations such as intersections can be modeled as vertices, while roads connecting them are represented as edges. By assigning weights to edges, additional information such as distance or expected travel time can be incorporated. 

One of the most important problems in graph theory is the shortest-path problem, which involves finding the minimum-cost path between two vertices. Classic algorithms for solving this problem, such as Dijkstra’s algorithm, are well known and widely taught in discrete mathematics and computer science courses \cite{cormen_leiserson_rivest_stein_2022} \cite{tardos_2006}. 

Although shortest-path algorithms are frequently used in routing applications, they rely on simplifying assumptions that may not hold in real transportation networks. Travel times on roads often vary due to congestion, traffic signals, or time-of-day effects. As a result, routes computed using fixed edge weights may not accurately reflect real travel behavior. This paper explores how shortest-path algorithms are used to model travel routes and examines the limitations of these models when applied to real-world transportation systems. 

\section{Mathematical Background}
A graph is defined as an ordered pair G=(V,E), where V is a set of vertices and E is a set of edges connecting pairs of vertices. In a weighted graph, each edge e exists in E is assigned a real-valued weight w(e), representing the cost of traversing that edge. 

In transportation models, vertices often correspond to physical locations, while edges represent direct travel connections such as roads. The weight of an edge may represent distance, average travel time, or another measure of cost. 

The shortest-path problem consists of finding a path between two vertices that minimizes the sum of the weights along that path. Dijkstra’s algorithm is a standard method for solving this problem in graphs with nonnegative edge weights \cite{dijkstra_1959}. The algorithm assumes that all edge weights are fixed and known in advance, and under these conditions it is guaranteed to find an optimal path. 

\section{Shortest Path Algorithms in Transportation Modeling}

Shortest-path algorithms are commonly used in routing and navigation systems. In these applications, transportation networks are modeled as graphs, and routes are selected based on criteria such as minimizing distance or travel time \cite{cormen_leiserson_rivest_stein_2022}. 

In practice, edge weights are often estimated using average or typical travel conditions. This approach can produce reasonable routes when traffic patterns are stable or when exact accuracy is not required. For example, in areas with low congestion or during off-peak hours, static shortest-path models may closely approximate actual travel times. 

However, this modeling approach simplifies the complexity of real transportation systems. Travel conditions can change significantly due to congestion, accidents, or traffic controls, which are not captured by fixed edge weights. As a result, shortest-path algorithms may produce routes that are optimal according to the model but suboptimal in practice. 

\bibliographystyle{plain} 
\bibliography{ref} 

\end{document}
